\input{estilo}
\begin{document}
\begin{figura}\begin{tikzpicture}
  % dos recuadros, el cliente y el servidor, cada uno con sus dos líneas de vida:
  % la aplicación y el stub. El tiempo corre hacia abajo.
  \foreach \x/\t in {-1.05/cliente, 4.65/servidor} {
    \draw[marco, draw=suave] (\x,0.55) rectangle ++(3.8,-6.0);
    \node[font=\small, above=2pt] at (\x+1.9,0.55) {\t};
  }
  \foreach \x/\t in {0/aplicación, 1.7/stub, 5.7/stub, 7.4/aplicación} {
    \draw[vida] (\x,0) -- (\x,-5.3);
    \node[etiq, above] at (\x,0) {\t};
  }
  \draw[msg] (0,-0.5) -- node[cod, font=\ttfamily\scriptsize, below, align=center] {add(200,\\id=7)} (1.7,-0.5);
  \foreach \k in {0,1,2} {
    \pgfmathsetmacro\y{-0.8-1.55*\k}
    \draw[msg] (1.7,\y) -- node[cod, font=\ttfamily\scriptsize, above, sloped, pos=0.5] {id=7} (5.7,\y-0.45);
    \draw[msg] (5.7,\y-0.45) -- (7.4,\y-0.45);
  }
  % las dos primeras respuestas se pierden; después de cada timeout el stub reintenta
  \foreach \k in {0,1} {
    \pgfmathsetmacro\y{-0.8-1.55*\k}
    \draw[respuesta] (5.7,\y-0.65) -- (3.7,\y-0.875);
    \pic[scale=0.6] at (3.7,\y-0.875) {rechazo};
    \draw (1.6,\y-1.4) -- (1.8,\y-1.4);
    \node[etiq, anchor=east] at (1.55,\y-1.4) {timeout};
  }
  % la tercera llega
  \draw[respuesta] (5.7,-4.55) -- (1.7,-5.0);
  \draw[respuesta] (1.7,-5.0) -- node[cod, above] {ok} (0,-5.0);
  % la primera vez se ejecuta y el id queda guardado en la base
  \node[evento=segundo] at (7.4,-1.25) {};
  \node[cod, text=segundo, anchor=north east] at (7.33,-1.32) {+200};
  \draw[msg, shorten <=3pt] (7.4,-1.25) -- (7.85,-1.25);
  \pic[scale=0.3] at (8.05,-1.42) {base};
  \node[cod, font=\ttfamily\scriptsize, below=1pt] at (8.05,-1.45) {ids};
  % las otras dos llegan con un id que ya está, y se ignoran
  \foreach \k in {1,2} {
    \pgfmathsetmacro\y{-0.8-1.55*\k}
    \pic[scale=0.5] at (7.4,\y-0.45) {rechazo};
  }
\end{tikzpicture}\end{figura}
\end{document}
